% GENERATED FILE. Edit canonical lessons, then run npm run generate:books.
% canonical-source-hash: fnv1a64:f288749170f63dc8
% canonical-lessons: IT-C500-cucire, IT-C500-ago, IT-C500-ricamare, IT-C500-uncinetto, IT-C500-rammendare

\chapter{Ago e filo}
\label{ch:it-ago-e-filo}
\hlchaptermodality{500}

\begin{chapteropening}
\textbf{By the end of this chapter:} \emph{I can talk about sewing, embroidery and mending.}

Complete the last lesson of chapter 500: I can talk about sewing, embroidery and mending.
\end{chapteropening}

\section[cucire]{cucire — to sew}

\label{lesson:IT-C500-cucire}



\begin{quote}
Before the new one: say the Italian for faithfulness, then the Italian for to confide.
\end{quote}

\subsection*{You'll want to know: cucire}
\textbf{cucire} — \textquotedblleft{}to sew\textquotedblright{}.

\begin{quote}
Mia nonna cuce i vestiti per tutta la famiglia.
\end{quote}

\begin{quote}
\emph{My grandmother sews clothes for the whole family.}
\end{quote}

Present: \emph{cucio}, \emph{cuci}, \emph{cuce}. \emph{La macchina da cucire} — the sewing machine.

\begin{cousinweb}
\emph{cucire} is inherited from Latin \emph{cōnsuere} \textquotedblleft{}to sew together\textquotedblright{} (Lewis \& Short), through a reshaped Vulgar Latin form whose details are not settled.
\end{cousinweb}

\begin{grammarlens}[title={What you've built}]
One more word for telling, explaining and describing at B1.
\end{grammarlens}

\subsection*{Your turn}
\begin{itemize}
\raggedright
  \item \emph{Say it:} \emph{cucire}
  \item \emph{Say it:} \emph{cucire}, once more
  \item \emph{Say it:} a sentence of your own with \emph{cucire}
  \item \emph{Recall:} say the Italian for faithfulness, then the Italian for to confide, then say \emph{cucire} again
\end{itemize}

\begin{tcolorbox}[breakable,colback=teal!4,colframe=teal!35!black,title={Before you move on}]
What does cucire mean? (\textquotedblleft{}To sew\textquotedblright{}.) Say it once more.
\end{tcolorbox}
\section[l'ago]{l'ago — needle}

\label{lesson:IT-C500-ago}



\begin{quote}
Before the new one: say the Italian for to confide, then the Italian for to sew.
\end{quote}

\subsection*{You'll want to know: l'ago}
\textbf{l'ago} — \textquotedblleft{}needle\textquotedblright{}.

\begin{quote}
Mi passi l'ago e il filo?
\end{quote}

\begin{quote}
\emph{Can you pass me the needle and thread?}
\end{quote}

The plural is \emph{gli aghi}. \emph{Cercare un ago in un pagliaio} — a needle in a haystack.

\begin{cousinweb}
\emph{ago} is inherited from Latin \emph{acus} \textquotedblleft{}a needle\textquotedblright{} (EtymDB; Lewis \& Short).
\end{cousinweb}

\begin{grammarlens}[title={What you've built}]
One more word for telling, explaining and describing at B1.
\end{grammarlens}

\subsection*{Your turn}
\begin{itemize}
\raggedright
  \item \emph{Say it:} \emph{l'ago}
  \item \emph{Say it:} \emph{l'ago}, once more
  \item \emph{Say it:} a sentence of your own with \emph{l'ago}
  \item \emph{Recall:} say the Italian for to confide, then the Italian for to sew, then say \emph{l'ago} again
\end{itemize}

\begin{tcolorbox}[breakable,colback=teal!4,colframe=teal!35!black,title={Before you move on}]
What does l'ago mean? (\textquotedblleft{}Needle\textquotedblright{}.) Say it once more.
\end{tcolorbox}
\section[ricamare]{ricamare — to embroider}

\label{lesson:IT-C500-ricamare}



\begin{quote}
Before the new one: say the Italian for to sew, then the Italian for needle.
\end{quote}

\subsection*{You'll want to know: ricamare}
\textbf{ricamare} — \textquotedblleft{}to embroider\textquotedblright{}.

\begin{quote}
Ha ricamato il suo nome sul fazzoletto.
\end{quote}

\begin{quote}
\emph{She embroidered her name on the handkerchief.}
\end{quote}

The noun is \emph{il ricamo}.

\begin{cousinweb}
\emph{ricamare} comes from Arabic \emph{raqama} \textquotedblleft{}to mark, to embroider\textquotedblright{} (EtymDB lists \emph{raqm} \textquotedblleft{}a mark\textquotedblright{}), probably through Medieval Latin.
\end{cousinweb}

\begin{grammarlens}[title={What you've built}]
One more word for telling, explaining and describing at B1.
\end{grammarlens}

\subsection*{Your turn}
\begin{itemize}
\raggedright
  \item \emph{Say it:} \emph{ricamare}
  \item \emph{Say it:} \emph{ricamare}, once more
  \item \emph{Say it:} a sentence of your own with \emph{ricamare}
  \item \emph{Recall:} say the Italian for to sew, then the Italian for needle, then say \emph{ricamare} again
\end{itemize}

\begin{tcolorbox}[breakable,colback=teal!4,colframe=teal!35!black,title={Before you move on}]
What does ricamare mean? (\textquotedblleft{}To embroider\textquotedblright{}.) Say it once more.
\end{tcolorbox}
\section[l'uncinetto]{l'uncinetto — crochet; crochet hook}

\label{lesson:IT-C500-uncinetto}



\begin{quote}
Before the new one: say the Italian for needle, then the Italian for to embroider.
\end{quote}

\subsection*{You'll want to know: l'uncinetto}
\textbf{l'uncinetto} — \textquotedblleft{}crochet; crochet hook\textquotedblright{}.

\begin{quote}
D'inverno faccio sciarpe all'uncinetto.
\end{quote}

\begin{quote}
\emph{In winter I crochet scarves.}
\end{quote}

\emph{Lavorare all'uncinetto} — to crochet; \emph{lavorare a maglia} — to knit.

\begin{cousinweb}
\emph{uncinetto} is a little form of \emph{uncino} \textquotedblleft{}hook\textquotedblright{}, inherited from Latin \emph{uncīnus} \textquotedblleft{}a hook\textquotedblright{} (Lewis \& Short), from \emph{uncus} \textquotedblleft{}bent\textquotedblright{}.
\end{cousinweb}

\begin{grammarlens}[title={What you've built}]
One more word for telling, explaining and describing at B1.
\end{grammarlens}

\subsection*{Your turn}
\begin{itemize}
\raggedright
  \item \emph{Say it:} \emph{l'uncinetto}
  \item \emph{Say it:} \emph{l'uncinetto}, once more
  \item \emph{Say it:} a sentence of your own with \emph{l'uncinetto}
  \item \emph{Recall:} say the Italian for needle, then the Italian for to embroider, then say \emph{l'uncinetto} again
\end{itemize}

\begin{tcolorbox}[breakable,colback=teal!4,colframe=teal!35!black,title={Before you move on}]
What does l'uncinetto mean? (\textquotedblleft{}Crochet; crochet hook\textquotedblright{}.) Say it once more.
\end{tcolorbox}
\section[rammendare]{rammendare — to darn, to mend}

\label{lesson:IT-C500-rammendare}



\begin{quote}
Before the new one: say the Italian for to embroider, then the Italian for crochet.
\end{quote}

\subsection*{You'll want to know: rammendare}
\textbf{rammendare} — \textquotedblleft{}to darn, to mend\textquotedblright{}.

\begin{quote}
Ho rammendato il buco nel calzino.
\end{quote}

\begin{quote}
\emph{I darned the hole in the sock.}
\end{quote}

For mending a hole by sewing; for other repairs say \emph{riparare}.

\begin{cousinweb}
\emph{rammendare} is built from \emph{r-} and \emph{ammendare} \textquotedblleft{}to correct\textquotedblright{}, from Latin \emph{ēmendāre} \textquotedblleft{}to free from faults\textquotedblright{} (Lewis \& Short), from \emph{menda} \textquotedblleft{}a fault\textquotedblright{}.
\end{cousinweb}

\begin{grammarlens}[title={What you've built}]
One more word for telling, explaining and describing at B1.
\end{grammarlens}

\subsection*{Your turn}
\begin{itemize}
\raggedright
  \item \emph{Say it:} \emph{rammendare}
  \item \emph{Say it:} \emph{rammendare}, once more
  \item \emph{Say it:} a sentence of your own with \emph{rammendare}
  \item \emph{Recall:} say the Italian for to embroider, then the Italian for crochet, then say \emph{rammendare} again
\end{itemize}

\begin{tcolorbox}[breakable,colback=teal!4,colframe=teal!35!black,title={Before you move on}]
What does rammendare mean? (\textquotedblleft{}To darn, to mend\textquotedblright{}.) Say it once more.
\end{tcolorbox}
